\documentclass[../../../../main.tex]{subfiles}
\begin{document}

\begin{flushright}
\footnotesize\textit{【自動文字起こし・要確認】}
\end{flushright}

[解]
(1) $n=1$では成立しているので、$n=k \in \mathbb{N}$での成立を仮定すると、
\begin{align*}
a_{k+1}^2 + b_{k+1}^2 &= 4a_k^2 + 4b_k^2 + 9c_k^2 - 2(4a_kb_k - 6a_kc_k - 6b_kc_k) \quad (\because \text{カテイ}) \\
&= c_{k+1}^2
\end{align*}
より$n=k+1$でも成立。以上から示された。■

(2) まず、$c_n > 0$を帰納的に示す。$n=1$で成立するので、$n=k \in \mathbb{N}$での成立を仮定する。
\begin{align*}
& c_{k+1} > 0 \Leftrightarrow 3c_k > 2(a_k+b_k) \Leftrightarrow 9c_k^2 > 4(a_k+b_k)^2 \quad (\because \text{両辺0以上}) \\
\Leftrightarrow\ & 9(a_k^2+b_k^2) > 4(a_k+b_k)^2 \Leftrightarrow 4(a_k-b_k)^2 + c_k^2 > 0
\end{align*}
だが、カテイより最後の不等式は成立するので、$n=k+1$でも成立。以上から$c_n > 0$である。■

次に、
\begin{align*}
& c_n > c_{n+1} \Leftrightarrow a_n+b_n > c_n \Leftrightarrow (a_n+b_n)^2 > c_n^2 \quad (\because \text{両辺0以上}) \\
\Leftrightarrow\ & 2a_nb_n > 0
\end{align*}
で、最後の不等式は成立するから、$c_n > c_{n+1}$である。■

(3) (2)から
\begin{equation*}
\begin{cases}
a_nb_n \neq 0 \implies c_n > c_{n+1} \\
a_nb_n = 0 \implies c_n = c_{n+1}
\end{cases}
\end{equation*}
だから、
\begin{equation*}
a_nb_n \neq 0, \quad a_{m+1}b_{m+1} = 0 \quad \cdots \text{\textcircled{1}}
\end{equation*}
である。以下 $a = a_m$などと表す。
\begin{itemize}
\item[$1^\circ\ a_{m+1} = 0\text{の時}$]
漸化式及び(1)から、
\begin{equation*}
\begin{cases} 2c = a+2b \\ c^2 = a^2+b^2 \end{cases}
\end{equation*}
$a$を消去して、
\begin{equation*}
(b-c)(5b-3c) = 0
\end{equation*}
$b=c$とすると $a^2+b^2=c^2$より $a=0$で、\textcircled{1}に反する。よって $5b = 3c$で、以上から
$a:b:c = 4:3:5$

\item[$2^\circ\ b_{m+1} = 0\text{の時}$]
$a_n, b_n$の対称性及び$1^\circ$から
$a:b:c = 3:4:5$
\end{itemize}
以上から
\begin{equation*}
a_m:b_m:c_m = 3:4:5 \text{ or } 4:3:5
\end{equation*}

\end{document}
